% Generated draft: requires mathematical and source review.
\documentclass[11pt]{article}
\usepackage[margin=25mm]{geometry}
\usepackage{fontspec}
\setmainfont{Latin Modern Roman}
\usepackage{amsmath,amssymb,mathtools,unicode-math}
\setmathfont{Latin Modern Math}
\usepackage{enumitem,needspace,xurl,hyperref}
\setlength{\parindent}{0pt}
\setlength{\parskip}{5pt}
\setlength{\emergencystretch}{3em}
\newcommand{\N}{\mathbb{N}}
\newcommand{\Z}{\mathbb{Z}}
\newcommand{\Q}{\mathbb{Q}}
\newcommand{\R}{\mathbb{R}}
\newcommand{\C}{\mathbb{C}}
\newcommand{\F}{\mathbb{F}}
\newcommand{\A}{\mathbb{A}}
\newcommand{\PP}{\mathbb{P}}
\newcommand{\E}{\mathbb{E}}
\newcommand{\eps}{\varepsilon}
\newcommand{\dd}{\,\mathrm{d}}
\begin{document}
\section*{UnsolvedMath ID: 20002294}
\textbf{Source problem code:} catalogue entry
\subsection{Definitions and mathematical statement}
A field $K$ is finitely generated if it is generated as a \emph{field} by finitely many elements over its prime field, $\Q$ or $\F_p$. Write $K\equiv L$ for elementary equivalence in the pure ring language without named generators or additional parameters, and $K\cong L$ for field isomorphism.

\textbf{Question 10, Florian Pop.} For finitely generated fields $K,L$, does
\[
 K\equiv L\quad\Longrightarrow\quad K\cong L?
\]
Both fields are required to be finitely generated; this is not categoricity among all fields or a decision procedure for their theories.

\textbf{Supplied Dittmann--Pop scope.} The aid reports single-sentence characterizations among finitely generated fields: for each $K$ in the stated range there is a sentence $\varphi_K$ such that a finitely generated $L$ satisfies it precisely when $L\cong K$. It consequently reports an affirmative answer unconditionally if $\operatorname{char}K\ne2$, and also in characteristic two when the absolute transcendence degree over $\F_2$ is at most three. In characteristic two and transcendence degree greater than three, it reports the conclusion only under the authors' stated resolution-of-singularities hypothesis above $\F_2$. The exact form of that hypothesis and a proving citation are not supplied; no unconditional resolution of this residual range is asserted. The characterizing sentence may depend on $K$; no effective uniform construction or decidability follows merely from its existence.

\textbf{Supplied positive-characteristic invariant package.} For a finitely generated field $K$ of characteristic $p>0$, put $d=\operatorname{trdeg}_{\F_p}K$ and $K^p=\{x^p:x\in K\}$. Then $[K:K^p]=p^d$. A first-order sentence expressing a $p$-basis of length $d$ is
\[
 \exists b_1\cdots b_d\ \forall x\ \exists!\,(c_e)_{e\in\{0,\ldots,p-1\}^d}
 \quad x=\sum_{e\in\{0,\ldots,p-1\}^d}
                       c_e^p\prod_{j=1}^db_j^{e_j}.
\]
For each fixed $p,d$ this is a finite formula; $d=0$ gives the perfect-field condition. The aid uses this to recover $d$ from the theory.

The relative algebraic closure of $\F_p$ in $K$ is a finite field $k_K=\F_{p^r}$ for some $r\ge1$. For every $n\ge1$, the number of roots in $K$ of $X^{p^n}-X$ is
\[
 p^{\gcd(r,n)}.
\]
Each finite root count is first-order expressible, and their full sequence determines $r$. Thus the aid's explicit elementary-equivalence consequences are
\[
 \operatorname{trdeg}_{\F_p}K=\operatorname{trdeg}_{\F_p}L,
 \qquad k_K\cong k_L.
\]
These invariants do not by themselves determine a general function field up to isomorphism and do not replace the scope-qualified characterization theorem.
\subsection{Short English statement}
Pop asks whether finitely generated fields with the same first-order theory must be isomorphic. The aid reports an affirmative characterization theorem outside a remaining conditional high-transcendence-degree characteristic-two range. Its explicit partial calculation recovers transcendence degree and the finite constant field in positive characteristic, without claiming those invariants alone determine the field.
\subsection{Sources}
\begin{itemize}
\item Local UnsolvedMath registry, permanent ID 20002294, source code catalogue entry; canonical \texttt{aim-logic-notes.json}, notes[69].
\item AIM \emph{Definability and decidability problems in number theory}, Question 10, attributed to Florian Pop: \url{https://aimath.org/pastworkshops/definabilityinntproblems.pdf}. The record reports agreement between PDF and JSON, not checked here.
\item Dittmann--Pop single-sentence characterization and its resolution-of-singularities qualification are named by the aid without full publication details, theorem label or exact resolution hypothesis.
\item Registry research attempt 1, dated 2026-08-14: partial result, rigor ``proved'', novelty confidence low. The explicitly machine-generated, non-peer-reviewed aid requires independent verification before citation.
\end{itemize}
\end{document}
